\paragraph{Learned representations and task estimators.}

A pretrained representation can improve an estimator without changing the object it estimates.
Several protein methods have the form
\begin{equation}
\boxed{
x\xrightarrow{\phi_\theta}z
\xrightarrow{\mathcal T_D}(\widehat y,\widehat\sigma)
\xrightarrow{\mathcal A}\text{ranking, retrieval, or experiment}.}
\label{eq:new-coordinates-old-engines}
\end{equation}
The encoder $\phi_\theta$ supplies learned coordinates, $\mathcal T_D$ fits task data, and
$\mathcal A$ selects an action.  In EVOLVEpro, a protein language model embeds sequences, a
random forest regresses measured activity, and successive experimental rounds supply new labels
\cite{jiang2025evolvepro}.  VIRAL instead uses an ESM3-derived representation with a
Gaussian-process surrogate and acquisition rule \cite{ito2025viral}.  PLMSearch uses learned
representations to retrieve candidates before sequence comparison supplies the homology ranking
\cite{liu2024plmsearch}.

These systems change the observation and estimator maps of the framework.  An assay prediction,
surrogate uncertainty, or retrieval score retains its own endpoint.  Extracting an interaction
operator would additionally require the source intervention, target readout, and reference
projection developed here.  Evaluation can locate the gain by fixing the representation while
comparing estimators, or fixing the estimator while comparing representations.  For an adaptive
experimental loop, the acquisition rule and newly measured data must also be included.
